\documentclass[11pt,a4paper]{article}
\usepackage[margin=27mm]{geometry}
\usepackage{amsmath,amsthm,amssymb,fontspec,unicode-math,hyperref}
\setmainfont{STIXGeneral}
\setmathfont{STIX Two Math}
\setmonofont{Menlo}
\hypersetup{hidelinks}
\usepackage{newunicodechar}
\newfontfamily\numeralfallback{STIX Two Text}
\newunicodechar{∎}{\ensuremath{\blacksquare}}
\newunicodechar{⁰}{{\numeralfallback\char"2070\relax}}
\newunicodechar{₀}{{\numeralfallback\char"2080\relax}}
\newunicodechar{₁}{{\numeralfallback\char"2081\relax}}
\newunicodechar{₂}{{\numeralfallback\char"2082\relax}}
\newunicodechar{₃}{{\numeralfallback\char"2083\relax}}
\newunicodechar{₄}{{\numeralfallback\char"2084\relax}}
\newunicodechar{⁻}{{\numeralfallback\char"207B\relax}}
\newunicodechar{⁸}{{\numeralfallback\char"2078\relax}}
\newunicodechar{ⁱ}{{\numeralfallback\char"2071\relax}}
\newunicodechar{⁴}{{\numeralfallback\char"2074\relax}}
\newunicodechar{⁷}{{\numeralfallback\char"2077\relax}}
\newunicodechar{⁶}{{\numeralfallback\char"2076\relax}}
\newunicodechar{⁵}{{\numeralfallback\char"2075\relax}}
\newunicodechar{₈}{{\numeralfallback\char"2088\relax}}
\newunicodechar{ᵢ}{\textsubscript{i}}
\newunicodechar{ⱼ}{\textsubscript{j}}
\newunicodechar{ₜ}{\textsubscript{t}}
\providecommand{\tightlist}{\setlength{\itemsep}{0pt}\setlength{\parskip}{0pt}}
\newtheorem{theorem}{Theorem}
\newtheorem{proposition}[theorem]{Proposition}
\newtheorem{lemma}[theorem]{Lemma}
\newtheorem{corollary}[theorem]{Corollary}
\theoremstyle{remark}\newtheorem{remark}[theorem]{Remark}
\newcommand{\gen}[1]{\langle #1\rangle}
\newcommand{\normalgen}[1]{\langle\!\langle #1\rangle\!\rangle}
\newcommand{\Z}{\mathbb Z}
\newcommand{\F}{\mathbb F}
\setlength{\emergencystretch}{2em}

\title{Square sections, alternating assemblies and free extensions}
\author{Human Compression Project}
\date{}
\begin{document}
\maketitle
\begin{abstract}
Compatible products give regular polytopes with any prescribed sequence of
two-power square sections. Coordinate kernels determine their exact orders,
flatness and deletion properties. A normal-kernel recognition theorem then
identifies their complete presentations, while jointly faithful deletions
identify the alternating two-branch construction. The same application also
recovers the three free-extension results through the torsion theorem for
amalgams. These are the six principal construction and extension theorems of
Cunningham, Feng, Hou and Schulte. The geometric assembly and free-extension
identifications are retained in explicit forms; their algebraic consequences
are proved here, including infinite inputs.
\end{abstract}

\section{Marked groups and geometric inputs}
A string C-group is generated by nonidentity involutions
$r_0,\ldots,r_{d-1}$, satisfies $[r_i,r_j]=1$ for $|i-j|>1$, and has
the intersection property
\[
 G_J\cap G_K=G_{J\cap K},\qquad G_J=\gen{r_j:j\in J}.
\]
We retain the regular-polytope correspondence stated precisely in [UA]:
such a marking gives a regular abstract polytope of rank $d$, with the
group acting regularly on flags. Its consecutive rank-three section on
labels $i,i+1,i+2$ has that marked subgroup as its group. A face of rank $j$
has stabilizer $G_{I\setminus\{j\}}$. Faces are incident when their
stabilizer cosets meet. These standard correspondences carry the original
attribution [MS02].

A deletion keeps a specified subset of the markers and sends all other
markers to the identity. UA.D proves that a family of all such deletions
gives the intersection property and
\[
 G=\normalgen{r_i:i\notin J}\rtimes G_J.
\]
At an initial or final cut, this split quotient is the flat amalgamation
property (FAP) used in [CFHS25]. We use this algebraic characterization.
Flatness is a different assertion: a polytope is $(j,k)$-flat if every
$j$-face is incident with every $k$-face. It will follow from an actual
factorization by face stabilizers, not merely from deletion maps.

Write $H_n$, $n\ge5$, for the affine square group in UA.Square. Its order
is $2^n$, its adjacent product orders are $4,4$, and its endpoint product
order is $2$. It has every Boolean deletion. Its complete affine model is
\[
 H_n=(\Z^2\rtimes D_8)/\Lambda_n,\qquad
 \Lambda_n=\begin{cases}
 2^e\Z^2,&n=2e+3,\\
 2^e\gen{(1,1),(1,-1)},&n=2e+4,
 \end{cases}\quad e\ge1.
\]
The full signed-permutation coordinate $D_8$ is retained even in a small
quotient on which its action is not faithful. The natural map $H_n\to H_5$
has kernel
\[
 N_n=\gamma_3(H_n)=2\Z^2/\Lambda_n,
 \qquad |N_n|=2^{n-5}.                                  \tag{1}
\]
It is injective on every proper marked subgroup: the three pairs retain
their exact orders and all singleton subgroups retain order two.

For completeness, UA.SquareRecognition proves that every finite two-group
string C-group of type $\{4,4\}$ is exactly one of these marked groups.
Its full proof uses orientation parity, the torsion-free translation kernel
of the square reflection presentation, and the Euclidean algorithm in
$\Z[i]$. Under the regular-polytope correspondence this gives the two
square-map families in the source's Lemma 4.1, one for each $n\ge5$.
Their group orders are $8s^2$ and $16s^2$, respectively, with $s=2^e$.
Thus the classification here follows from a complete shared proof, rather
than an unproved exhaustion of the affine examples. The finite construction
below itself needs only the explicit groups $H_n$.

\section{One construction for every sequence of square sections}
Let $T$ be a finite forest with vertex set $I$. The coordinate group
$B_T$ of UA.Forest has binary coordinates on the vertices and edges. After
choosing an orientation for each edge, multiplication is
\[
 (v,z)(w,t)=(v+w,z+t+\beta(v,w)),\qquad
 \beta_{i\to j}(v,w)=v_jw_i.
\]
Its order is $2^{|I|+|E(T)|}$; deleting coordinates gives every Boolean
deletion. Each two-edge path has marked group $H_5$.
Choose distinct two-edge paths $S$ and parameters $n_S\ge5$. Attach
$H_{n_S}\to (B_T)_S=H_5$ on each support. These maps are faithful on proper
marked subgroups by (1), and the supports form an antichain. UA.SK and its
local, deletion and lower-central clauses give the complete compatible
product $G$, with its actual marker tuples, and
\begin{align}
 |G_J|&=2^{|J|+|E(T[J])|+\sum_{S\subseteq J}(n_S-5)},       \tag{2}\\
 \gamma_3(G)&=\prod_S N_{n_S},\qquad G/\gamma_3(G)=B_T.     \tag{3}
\end{align}
It has all Boolean deletions and exact local groups $G_S=H_{n_S}$.
The order in (2) is the order of the group generated by the marker tuples:
generation of the whole compatible product is part of UA.SK.

\begin{theorem}[Cunningham--Feng--Hou--Schulte, Theorem 5.1]
Let $d\ge3$ and $n_3,\ldots,n_d\ge5$. There is a regular rank-$d$
polytope $\mathcal P$ whose consecutive square sections have groups
$H_{n_3},\ldots,H_{n_d}$, with
\[
 |\Gamma(\mathcal P)|=2^{\sum_{j=3}^{d} n_j-3(d-3)}.        \tag{4}
\]
For $d\ge4$ it is $(d-4,d-1)$-flat, hence flat. It has the FAP with
respect to its co-$(d-4)$-faces, co-$(d-3)$-faces and facets.
At $d=3$ the co-$(-1)$-face cut means the whole polytope.
\end{theorem}
\begin{proof}
Take the path $0,1,\ldots,d-1$ and attach $H_{n_{i+3}}$ on every
$S_i=\{i,i+1,i+2\}$, $0\le i\le d-3$. Distant markers commute in
the base and in every local factor, so the resulting IP marking is string.
Its generators have exact order two by their base images. Local recognition
gives the prescribed marked square sections. Equation (2), with $J=I$,
gives
\[
 |G|=2^{2d-1+\sum_{j=3}^{d}(n_j-5)}
     =2^{\sum_{j=3}^{d}n_j-3(d-3)}.
\]
The regular-polytope correspondence now gives the stated object and order.

All deletions exist. In particular delete onto
$\{d-3,d-2,d-1\}$, onto $\{d-2,d-1\}$, and onto
$\{0,\ldots,d-2\}$. UA.D supplies the three split normal-closure
decompositions, which are exactly the stated FAP assertions. For $d=3$
the first deletion is the identity, so its interpretation needs no exception
to the algebraic statement.

For flatness assume $d\ge4$, and put
$L=\{0,\ldots,d-2\}$ and $M=\{d-3,d-2,d-1\}$. Every support is
contained in $L$ or $M$. The vertex and edge coordinates give
$B_T=(B_T)_L(B_T)_M$: one can first choose the two factors' vertex
coordinates, then assign each central edge coordinate to a side containing
that edge, correcting the product's bilinear contribution there.
UA.SK.Separator gives $G=G_LG_M$. The first factor is a facet stabilizer;
the second lies in the stabilizer of the base $(d-4)$-face. Thus their
full stabilizers multiply to $G$. Any right coset of one meets any right
coset of the other, so every such face is incident with every facet.
Every vertex lies in a $(d-4)$-face; it is therefore incident with every
facet as well. This proves the two flatness conclusions.
\end{proof}

The group is also precisely the particular iterated mix used in the source,
not merely another group with its order and sections. Projection into
$\prod_i H_{n_{i+3}}$ is injective: an element with all local coordinates
trivial has every triple deletion of its base coordinate trivial, and the
triples cover every vertex and edge coordinate of $B_T$. Its image is
generated by the marked tuples in those factors. Equation (8) of the source
and its induction in Theorem 5.1 produce exactly these tuples, with only
reassociation of the direct product. The marked groups are consequently
isomorphic.

\section{Recognizing the complete universal presentation}
Let $P_0$ have the string Coxeter relations of type $\{4,\ldots,4\}$
and the complete lattice relations of $H_{n_{i+3}}$ on every consecutive
triple. Let $P$ be its quotient by
\[
 [(r_i r_{i+1})^2,(r_{i+2}r_{i+3})^2]=1
       \quad(0\le i\le d-4).                             \tag{5}
\]
The following local identity supplies the required normalization; its proof
is common to every rank and parameter.

\begin{lemma}
If involutions $a,b,c,d$ satisfy $[a,c]=[a,d]=[b,d]=1$ and
$[(ab)^2,(cd)^2]=1$, then $d$ commutes with both $(cbab)^2$ and
$(abcb)^2$.
\end{lemma}
\begin{proof}
Put $x=bab$, $A=(ab)^2=ax$, and $F=(cd)^2$. Since $a$ commutes with
$F$, the relation $[A,F]=1$ gives $[x,F]=1$. Also $[x,d]=1$, hence
$[x,Fd]=[x,cdc]=1$. Conjugating by $c$ gives $[cxc,d]=1$.
Both factors of $(cbab)^2=(cxc)x$ commute with $d$. Finally
$abcb=b(cbab)^{-1}b$, and $b$ commutes with $d$.
\end{proof}

\begin{theorem}[Cunningham--Feng--Hou--Schulte, Theorem 5.3]
The group $P_0$ is the universal regular-polytope group with the indicated
consecutive square sections. The group constructed in Theorem 1 is precisely
its marked quotient $P$ by (5).
\end{theorem}
\begin{proof}
Let $H$ denote the group of Theorem 1. Its local sections have the exact
complete presentations used in $P_0$. Relations (5) hold in its class-two
base. In each local triple coordinate at least one of the two disjoint
edges loses an endpoint, so that pair product has square one. Thus every
relator of $P$ holds in $H$, giving a marked surjection $P\to H$.
Each local square group maps to $P$ by the defining relations; projection
onward to its own coordinate of $H$ is the identity, so this local map is
faithful.

Let $K_i$ be the image of $\gamma_3(H_{n_{i+3}})$ in $P$. It is generated
by the two doubled translations appearing in the lemma. The three local
markers normalize it. The lemma and its reversed version show that the
nearest outside markers centralize its generators; every farther marker
commutes with all local markers by the string relations. Hence each
$K_i$ is normal in $P$.

Kill all the $K_i$. Every consecutive triple becomes class two. Every
adjacent square therefore commutes with its endpoints and its nearest
outside markers; the string relations give commutation with all remaining
markers. The quotient is class two and satisfies the complete path-base
presentation of UA.Forest. Conversely its map to $H$ followed by the base
map gives the marked surjection onto that same path base. The two maps are
inverse on generators, so this is the exact residual quotient. All four
hypotheses of UA.Recognition now hold. It identifies $P$ with $H$ without
assuming $P$ finite or estimating its order.

Finally every defining relator of $P_0$ survives all deletion assignments.
Even Coxeter relators do so directly; UA.Square proves the assertion for
each complete local lattice relation. Thus $P_0$ has IP by UA.D. Its
generators are nonidentity involutions, and each local group embeds because
its map onward to $H$ is faithful. Hence $P_0$ is a string C-group with
the required exact sections. Every other such marked group satisfies all
of its relators and is its marked quotient. The standard covering
correspondence gives the universal regular polytope. No finiteness of
$P_0$ is used or asserted.
\end{proof}

In the smallest diagonal square group $H_6$, the two doubled translations
coincide. Formula (1) uses their full lattice subgroup and does not assume
that their two cyclic subgroups are independent. This preserves the
previously recorded qualification to the source's presentation proof.

\section{Two branches and the full alternating automorphism group}
We retain the following precise FAP case of the alternating-assembly theorem
of Monson and Schulte [MSA, Theorem 5.3]. If regular $(d-1)$-polytopes
$\mathcal P,\mathcal Q$ have an identified common facet $\mathcal K$ and
FAP decompositions
\[
 \Gamma(\mathcal P)=N_P\rtimes\Gamma(\mathcal K),\qquad
 \Gamma(\mathcal Q)=N_Q\rtimes\Gamma(\mathcal K),
\]
then the tail-triangle group
$(N_P\times N_Q)\rtimes\Gamma(\mathcal K)$, with the two distinct final
markers, gives an alternating semiregular $d$-polytope with those facets.
Its full automorphism group is that group when the two facets are not
isomorphic, and its extension by the branch-swapping $C_2$ when they are
the same marked polytope. In the latter case the polytope is regular.
This geometric and full-automorphism identification is retained; group IP
alone would not establish it.

\begin{theorem}[Cunningham--Feng--Hou--Schulte, Theorem 6.1]
For $d\ge4$, $n_3,\ldots,n_{d-1}\ge5$ and $m_{d-1}\ge5$, there is an
alternating semiregular rank-$d$ polytope with regular facet lists
\[
 (n_3,\ldots,n_{d-2},n_{d-1}),\qquad
 (n_3,\ldots,n_{d-2},m_{d-1}).
\]
Its full automorphism group has order
\[
 2^{\sum_{j=3}^{d-2}n_j+n_{d-1}+m_{d-1}-3(d-3)}             \tag{6}
\]
when the last parameters differ, and twice this order when they agree.
For $d=4$ the shared list is empty and the common facet is a square.
\end{theorem}
\begin{proof}
Take the tree with common prefix $0,\ldots,d-3$ and two final leaves
$u,v$, both adjacent to $d-3$. Attach the two prescribed lists of square
groups along the two paths, with the shared prefix supports counted once.
The construction leading to (2) applies. Let $L$ and $M$ be the two full
path label sets, and put $C=L\cap M$.
Local recognition and the coordinate construction identify $G_L,G_M$ with
the two groups of Theorem 1. Their common subgroup $G_C$ is its shorter
path group, or the square dihedral group of order eight when $d=4$.

The deletions onto $L$ and $M$ are jointly faithful. Every vertex and edge
coordinate of the base is retained on at least one side, and every local
kernel support lies entirely on a side, where its coordinate is retained.
An element killed by both deletions therefore has every coordinate trivial.
UA.Mix identifies the group with the full compatible product
$G_L\times_{G_C}G_M$. This step uses joint coordinate detection and actual
generation of the compatible product, not an order comparison.
By UA.Fiber it is exactly
$(\ker(G_L\to G_C)\times\ker(G_M\to G_C))\rtimes G_C$, with the original
tail-triangle marking. The retained assembly theorem applies.

The common group has order
$2^{\sum_{j=3}^{d-2}n_j-3(d-5)}$, including $8$ for the empty list at
$d=4$. The two kernels have orders $2^{n_{d-1}-3}$ and
$2^{m_{d-1}-3}$ by (2). Multiplication gives (6). Unequal last parameters
give unequal facet-group orders, so their polytopes cannot be isomorphic.
Equal parameters give the identical marked construction and its branch
interchange. The full-group alternative in the retained theorem supplies
exactly the factor of two asserted.
\end{proof}

The co-$(d-4)$-section not lying inside one of the two facet lists is not
separately prescribed by this theorem. No extra statement about that section
is used to prove (6).

\section{Free extensions and two-primary torsion}
Following the source, an infinite two-group means a finitely generated
infinite group every finite-order element of which has two-power order.
It need not be a torsion group. UA.Amalgam proves that this torsion property
passes to subgroups, to extensions with finite two-group quotient, and to
amalgamated free products. Its proof retains the precise amalgam normal-form
theorem and proves infinitude when the common subgroup is proper in both
factors.

We retain three geometric identifications, including existence, universality,
rank and symmetry type:
\begin{enumerate}
\item A regular rank-$d$ polytope $\mathcal K$ with facet $\mathcal F$ has
a free regular extension of rank $d+1$, with group
\[
 \Gamma(\mathcal K)*_{\Gamma(\mathcal F)}
       (\Gamma(\mathcal F)\times C_2).
\]
This is the standard regular free-extension construction [MS02,4D4].
\item An oriented chiral rank-$d$ polytope $\mathcal K$ with directly
regular facets $\mathcal F$ has a free oriented chiral extension of rank
$d+1$, with group
\[
 \Gamma(\mathcal K)*_{\Gamma^+(\mathcal F)}\Gamma(\mathcal F),
 \qquad \Gamma(\mathcal F)=\Gamma^+(\mathcal F)\rtimes C_2.
\]
Here the facet rotation subgroup embeds in $\Gamma(\mathcal K)$.
This is the free chiral-extension theorem of Schulte and Weiss [SW95].
\item Regular rank-$d$ polytopes $\mathcal P,\mathcal Q$ with isomorphic
facets $\mathcal K$ have a universal alternating semiregular extension of
rank $d+1$. Its group is
\[
 \Gamma(\mathcal P)*_{\Gamma(\mathcal K)}\Gamma(\mathcal Q)
\]
when $\mathcal P\not\cong\mathcal Q$, and the extension of this group by
the branch-swapping $C_2$ when $\mathcal P\cong\mathcal Q$. This is the
universal alternating construction of Monson and Schulte [MST12,MST21].
\end{enumerate}
All input ranks in these statements satisfy the source's convention $d\ge3$.
The amalgamating subgroup is proper in each factor: in the polytope factor
it is a proper face or facet-rotation stabilizer, while in the added factor
of the first two constructions it has index two. These are the actual
geometric subgroup embeddings, not abstract groups of the same order.

\begin{theorem}[Cunningham--Feng--Hou--Schulte, Theorems 7.1--7.3]
In each of the three identifications above, if the input automorphism groups
are finite or infinite two-groups in the stated convention, the full
automorphism group of the resulting extension is a finitely generated
infinite group with only two-power finite torsion.
\end{theorem}
\begin{proof}
For the regular extension, $\Gamma(\mathcal F)$ is a subgroup of the input
group and has the torsion property. Its product with $C_2$ has the property
by the finite-extension clause of UA.Amalgam. For the chiral extension it
is the rotation subgroup $\Gamma^+(\mathcal F)$ that embeds in the input
group. Its index-two extension $\Gamma(\mathcal F)$ again has the
property. This does not assume that the full facet group embeds in the
chiral automorphism group. In the alternating case both factors are input
groups and already have the required property.

In each case every finite-order element of the amalgam is conjugate into
one of its factors, so it has two-power order. Properness on both sides
provides reduced alternating words of arbitrarily large length and proves
that the amalgam is infinite. All factors are finitely generated by their
finite marked lists, so their amalgam is finitely generated. In the equal
alternating-facet case, its additional extension by $C_2$ still has the
torsion property and remains infinite and finitely generated. The retained
geometric identifications supply the rank, symmetry and universal-extension
conclusions in the three source theorems.
\end{proof}

\section{Coverage and attribution}
The source edition is arXiv:2512.15511v1, dated 18 December 2025. We have
proved its Theorem 5.1(a--d), complete Theorem 5.3, Theorem 6.1(a--b), and
Theorems 7.1--7.3, relative to the explicitly retained geometric statements.
The full square-recognition theorem in the umbrella also recovers
Lemma 4.1. The source's unnumbered remarks and power-polytope
discussion are not included in this principal contract.

The finite constructions use the same supported-kernel theorem. Its
normal-kernel recognition clause proves the complete presentation, and
joint coordinate detection proves the full two-branch product. Local
normalization and the base quotient remain explicit work. The infinite
results use a different common mechanism, the amalgam normal form; they
are not inferred from the finite construction. The original results and
all earlier geometric inputs retain their stated attribution.

\begin{thebibliography}{9}
\bibitem{UA}[UA] \emph{Supported kernels and finite two-group quotients},
companion umbrella in this volume.
\bibitem{CFHS25}[CFHS25] G. Cunningham, Y.-Q. Feng, D.-D. Hou and E. Schulte,
\emph{Symmetric polytopes whose automorphism groups are 2-groups},
arXiv:2512.15511v1, 18 December 2025.
\bibitem{MS02}[MS02] P. McMullen and E. Schulte,
\emph{Abstract Regular Polytopes}, Cambridge University Press, 2002.
\bibitem{MSA}[MSA] B. Monson and E. Schulte,
\emph{The assembly problem for alternating semiregular polytopes},
\emph{Discrete \& Computational Geometry} 64 (2020), 453--482,
Theorem 5.3, in the precise FAP form retained above.
\bibitem{SW95}[SW95] E. Schulte and A. I. Weiss,
\emph{Free extensions of chiral polytopes}, \emph{Canadian Journal of
Mathematics} 47(3) (1995), 641--654.
\bibitem{MST12}[MST12] B. Monson and E. Schulte,
\emph{Semiregular polytopes and amalgamated C-groups},
\emph{Advances in Mathematics} 229 (2012), 2767--2791.
\bibitem{MST21}[MST21] B. Monson and E. Schulte,
\emph{Universal alternating semiregular polytopes},
\emph{Canadian Journal of Mathematics} 73 (2021), 572--595.
The exact group and full-automorphism interface is stated above.
\end{thebibliography}
\end{document}
